\documentclass{article}
\usepackage{mathtools} 
\usepackage{fontspec}
\usepackage[UTF8]{ctex}
\usepackage{amsthm}
\usepackage{mdframed}
\usepackage{xcolor}
\usepackage{amssymb}
\usepackage{amsmath}
\usepackage{tikz}
\usepackage{pgfplots}
\usepgfplotslibrary{polar}
% \pgfplotsset{compat=1.18}


% 定义新的带灰色背景的说明环境 zremark
\newmdtheoremenv[
  backgroundcolor=gray!10,
  % 边框与背景一致，边框线会消失
  linecolor=gray!10
]{zremark}{说明}

% 通用矩阵命令: \flexmatrix{矩阵名}{元素符号}{行数}{列数}
\newcommand{\flexmatrix}[4]{
  \[
  #1 = \begin{pmatrix}
    #2_{11}     & #2_{12}     & \cdots & #2_{1#4}   \\
    #2_{21}     & #2_{22}     & \cdots & #2_{2#4}   \\
    \vdots      & \vdots      & \ddots & \vdots     \\
    #2_{#31}    & #2_{#32}    & \cdots & #2_{#3#4}
  \end{pmatrix}
  \]
}

% 简化版命令（默认矩阵名为A，元素符号为a）: \quickmatrix{行数}{列数}
\newcommand{\quickmatrix}[2]{\flexmatrix{A}{a}{#1}{#2}}

\begin{document}
\title{注释 11.2}
\author{张志聪}
\maketitle

\begin{zremark}
  $\int_{a}^{+\infty} f(x) dx = \int_{a}^{+\infty} g(x) dx + \int_{a}^{+\infty} f(x) dx$，
  只有当证明第一个积分发散（收敛），第二个积分收敛（发散），才能说明
  $\int_{a}^{+\infty} f(x) dx$发散。
\end{zremark}

\textbf{证明：}

反例：一个收敛的无穷积分，可以拆为两个发散的无穷积分。比如
\begin{align*}
  \int_{a}^{+\infty} \frac{1}{x^2} dx
   & = \int_{a}^{+\infty} \frac{x}{x^2} - \frac{x - 1}{x^2} dx \\
   & = \int_{a}^{+\infty} \frac{1}{x} - \frac{x - 1}{x^2} dx
\end{align*}

注：本身这种无穷积分拆分成两个无穷积分的和，其实是利用了极限的四则运算法则。

\end{document}